% GENERATED FILE. Edit canonical lessons, then run npm run generate:books.
% canonical-source-hash: fnv1a64:a3052fa9127cb4f2
% canonical-lessons: JA-C237-shashin, JA-C237-toru, JA-C237-atsumeru, JA-C237-atsumaru, JA-C237-ningyou

\chapter{Photos and Collections}
\label{ch:ja-photos-and-collections}
\hlchaptermodality{237}

\begin{chapteropening}
\textbf{By the end of this chapter:} \emph{I can talk about photos I took and things I collected.}

Complete the last lesson of chapter 237: I can talk about photos I took and things I collected.
\end{chapteropening}

\section[\texorpdfstring{\ja{しゃしん} (shashin)}{しゃしん (shashin)}]{\ja{しゃしん} (shashin) — a photo; a photograph}

\label{lesson:JA-C237-shashin}



\begin{quote}
Before the new one: say the Japanese for music, then the Japanese for a musical performance.
\end{quote}

\subsection*{You'll want to know: \ja{しゃしん}}
\textbf{\ja{しゃしん}} — \emph{shashin} — \textquotedblleft{}a photo; a photograph\textquotedblright{}.

\begin{quote}
\textbf{\ja{きのう} \ja{かぞくの} \ja{しゃしんを} \ja{みました。}} — \emph{kinō kazoku no shashin o mimashita} — \textquotedblleft{}Yesterday I looked through family photos.\textquotedblright{}
\end{quote}

\begin{cousinweb}
\textbf{\ja{しゃしん}} (\emph{shashin}) is Sino-Japanese (kango): two Chinese characters read with their on'yomi, readings borrowed from Chinese. They are \emph{sha}, \textquotedblleft{}copy\textquotedblright{} (Middle Chinese \emph{sjaeX}) and \emph{shin}, \textquotedblleft{}true\textquotedblright{} (\emph{tsyin}) (KANJIDIC2; Middle Chinese in Baxter's transcription).
\end{cousinweb}

\begin{grammarlens}[title={What you've built}]
One more word for everyday life at A2.
\end{grammarlens}

\subsection*{Your turn}
\begin{itemize}
\raggedright
  \item \emph{Say it:} \emph{shashin}
  \item \emph{Say it:} \emph{shashin}, once more
  \item \emph{Say it:} \emph{kinō kazoku no shashin o mimashita} — Yesterday I looked through family photos.
  \item \emph{Recall:} say the Japanese for music, then the Japanese for a musical performance, then say \emph{shashin} again
\end{itemize}

\begin{tcolorbox}[breakable,colback=teal!4,colframe=teal!35!black,title={Before you move on}]
What does \ja{しゃしん} mean? (\textquotedblleft{}A photo; a photograph\textquotedblright{}.) Say it once more.
\end{tcolorbox}
\section[\texorpdfstring{\ja{とる} (toru)}{とる (toru)}]{\ja{とる} (toru) — to photograph; to take (a photo)}

\label{lesson:JA-C237-toru}



\begin{quote}
Before the new one: say the Japanese for a musical performance, then the Japanese for a photo.
\end{quote}

\subsection*{You'll want to know: \ja{とる}}
\textbf{\ja{とる}} — \emph{toru} — \textquotedblleft{}to photograph; to take (a photo)\textquotedblright{}.

\begin{quote}
\textbf{\ja{せんげつ} \ja{やまで} \ja{しゃしんを} \ja{とりました。}} — \emph{sengetsu yama de shashin o torimashita} — \textquotedblleft{}Last month I took photos in the mountains.\textquotedblright{}
\end{quote}

Polite \textbf{\ja{とります}}, \emph{torimasu}; a godan verb.

\begin{cousinweb}
\textbf{\ja{とる}} (\emph{toru}) is a native Japanese word (wago), inherited rather than borrowed from Chinese; its kanji was attached later as a kun'yomi (JMdict, KANJIDIC2).
\end{cousinweb}

\begin{grammarlens}[title={What you've built}]
One more word for everyday life at A2.
\end{grammarlens}

\subsection*{Your turn}
\begin{itemize}
\raggedright
  \item \emph{Say it:} \emph{toru}
  \item \emph{Say it:} \emph{toru}, once more
  \item \emph{Say it:} \emph{sengetsu yama de shashin o torimashita} — Last month I took photos in the mountains.
  \item \emph{Recall:} say the Japanese for a musical performance, then the Japanese for a photo, then say \emph{toru} again
\end{itemize}

\begin{tcolorbox}[breakable,colback=teal!4,colframe=teal!35!black,title={Before you move on}]
What does \ja{とる} mean? (\textquotedblleft{}To photograph; to take (a photo)\textquotedblright{}.) Say it once more.
\end{tcolorbox}
\section[\texorpdfstring{\ja{あつめる} (atsumeru)}{あつめる (atsumeru)}]{\ja{あつめる} (atsumeru) — to collect}

\label{lesson:JA-C237-atsumeru}



\begin{quote}
Before the new one: say the Japanese for a photo, then the Japanese for to photograph.
\end{quote}

\subsection*{You'll want to know: \ja{あつめる}}
\textbf{\ja{あつめる}} — \emph{atsumeru} — \textquotedblleft{}to collect\textquotedblright{}.

\begin{quote}
\textbf{\ja{きょねん} \ja{おとうとは} \ja{きっぷを} \ja{あつめました。}} — \emph{kyonen otōto wa kippu o atsumemashita} — \textquotedblleft{}Last year my younger brother collected tickets.\textquotedblright{}
\end{quote}

Polite \textbf{\ja{あつめます}}, \emph{atsumemasu}; an ichidan verb.

\begin{cousinweb}
\textbf{\ja{あつめる}} (\emph{atsumeru}) is a native Japanese word (wago), inherited rather than borrowed from Chinese; its kanji was attached later as a kun'yomi (JMdict, KANJIDIC2).
\end{cousinweb}

\begin{grammarlens}[title={What you've built}]
One more word for everyday life at A2.
\end{grammarlens}

\subsection*{Your turn}
\begin{itemize}
\raggedright
  \item \emph{Say it:} \emph{atsumeru}
  \item \emph{Say it:} \emph{atsumeru}, once more
  \item \emph{Say it:} \emph{kyonen otōto wa kippu o atsumemashita} — Last year my younger brother collected tickets.
  \item \emph{Recall:} say the Japanese for a photo, then the Japanese for to photograph, then say \emph{atsumeru} again
\end{itemize}

\begin{tcolorbox}[breakable,colback=teal!4,colframe=teal!35!black,title={Before you move on}]
What does \ja{あつめる} mean? (\textquotedblleft{}To collect\textquotedblright{}.) Say it once more.
\end{tcolorbox}
\section[\texorpdfstring{\ja{あつまる} (atsumaru)}{あつまる (atsumaru)}]{\ja{あつまる} (atsumaru) — to gather; to come together}

\label{lesson:JA-C237-atsumaru}



\begin{quote}
Before the new one: say the Japanese for to photograph, then the Japanese for to collect.
\end{quote}

\subsection*{You'll want to know: \ja{あつまる}}
\textbf{\ja{あつまる}} — \emph{atsumaru} — \textquotedblleft{}to gather; to come together\textquotedblright{}.

\begin{quote}
\textbf{\ja{しゅうまつに} \ja{ともだちが} \ja{うちに} \ja{あつまりました。}} — \emph{shūmatsu ni tomodachi ga uchi ni atsumarimashita} — \textquotedblleft{}At the weekend my friends gathered at my house.\textquotedblright{}
\end{quote}

Polite \textbf{\ja{あつまります}}, \emph{atsumarimasu}; a godan verb. Its partner \textbf{\ja{あつめる}}, \emph{atsumeru}, takes an object.

\begin{cousinweb}
\textbf{\ja{あつまる}} (\emph{atsumaru}) is a native Japanese word (wago), inherited rather than borrowed from Chinese; its kanji was attached later as a kun'yomi (JMdict, KANJIDIC2).
\end{cousinweb}

\begin{grammarlens}[title={What you've built}]
One more word for everyday life at A2.
\end{grammarlens}

\subsection*{Your turn}
\begin{itemize}
\raggedright
  \item \emph{Say it:} \emph{atsumaru}
  \item \emph{Say it:} \emph{atsumaru}, once more
  \item \emph{Say it:} \emph{shūmatsu ni tomodachi ga uchi ni atsumarimashita} — At the weekend my friends gathered at my house.
  \item \emph{Recall:} say the Japanese for to photograph, then the Japanese for to collect, then say \emph{atsumaru} again
\end{itemize}

\begin{tcolorbox}[breakable,colback=teal!4,colframe=teal!35!black,title={Before you move on}]
What does \ja{あつまる} mean? (\textquotedblleft{}To gather; to come together\textquotedblright{}.) Say it once more.
\end{tcolorbox}
\section[\texorpdfstring{\ja{にんぎょう} (ningyō)}{にんぎょう (ningyō)}]{\ja{にんぎょう} (ningyō) — a doll; a puppet}

\label{lesson:JA-C237-ningyou}



\begin{quote}
Before the new one: say the Japanese for to collect, then the Japanese for to gather.
\end{quote}

\subsection*{You'll want to know: \ja{にんぎょう}}
\textbf{\ja{にんぎょう}} — \emph{ningyō} — \textquotedblleft{}a doll; a puppet\textquotedblright{}.

\begin{quote}
\textbf{\ja{いもうとは} \ja{きょねん} \ja{にんぎょうを} \ja{あつめました。}} — \emph{imōto wa kyonen ningyō o atsumemashita} — \textquotedblleft{}Last year my younger sister collected dolls.\textquotedblright{}
\end{quote}

\begin{cousinweb}
\textbf{\ja{にんぎょう}} (\emph{ningyō}) is Sino-Japanese (kango): two Chinese characters read with their on'yomi, readings borrowed from Chinese. They are \emph{nin}, \textquotedblleft{}person\textquotedblright{} (Middle Chinese \emph{nyin}) and \emph{gyō}, \textquotedblleft{}shape\textquotedblright{} (\emph{heng}) (KANJIDIC2; Middle Chinese in Baxter's transcription).
\end{cousinweb}

\begin{grammarlens}[title={What you've built}]
One more word for everyday life at A2.
\end{grammarlens}

\subsection*{Your turn}
\begin{itemize}
\raggedright
  \item \emph{Say it:} \emph{ningyō}
  \item \emph{Say it:} \emph{ningyō}, once more
  \item \emph{Say it:} \emph{imōto wa kyonen ningyō o atsumemashita} — Last year my younger sister collected dolls.
  \item \emph{Recall:} say the Japanese for to collect, then the Japanese for to gather, then say \emph{ningyō} again
\end{itemize}

\begin{tcolorbox}[breakable,colback=teal!4,colframe=teal!35!black,title={Before you move on}]
What does \ja{にんぎょう} mean? (\textquotedblleft{}A doll; a puppet\textquotedblright{}.) Say it once more.
\end{tcolorbox}
